She had tried four things since the arm, and none of them had held.

She had tried the water lines first, because Roland Storm asked her plainly whether she wanted work and she had said yes before she’d thought about what the work would actually be, and it turned out to be standing in a street telling frightened people where the line ended, which she was good at and could not do for more than a morning at a stretch before the wanting-to-leave got too loud to argue with.

She had tried the Unsullied drilling grounds, once, at Grey Worm’s own invitation, on the theory that a woman who could still hold a sword might have something to teach men relearning how to fight without the other nine of them beside them. She had lasted two sessions. She was good at it. She had also discovered, on the second morning, that she did not want to be good at teaching men to kill more efficiently, not any longer, and had not been able to say why to Grey Worm in any way that made sense to either of them, and had simply stopped coming.

She had tried, briefly and with a kind of grim curiosity, sitting in on one of Sam’s record afternoons, to see whether the keeping of things down on paper was a trade a person with one working hand might take up. She had lasted an hour. Writing, it turned out, required a great deal more patience than killing ever had.

The fourth thing she had tried was Gendry, and that one had held longer than any of the others, and she had known for some while now, in the flat unhurried way she knew most true things about herself, that it was going to be the hardest one to put down.

\sceneBreak

The forge caught on the ordinary afternoon that ended it, which she thought about a great deal afterward and never found a way to make sound as small as it had actually been.

He had two boys working the bellows now, Flea Bottom lads of about twelve who had never held a hammer before the war and were slowly becoming smiths the way anyone becomes anything, by being needed at it, and one of them knocked a coal basket wrong and a heap of rag and shaving caught at the base of the second forge before anybody had properly registered that it had happened.

Gendry moved before Arya had finished turning her head toward the sound.

He had the basket over, the flame smothered under a wet sack, and the boy — who had gone white and useless the way any boy goes white and useless the first time he causes a thing like that — bodily hauled two steps back from the heat, all inside about the length of time it took Arya to get up off the bench she’d been sitting on. It was over before it was properly a story. Nobody outside that yard would ever hear of it. He set the sack down, said something short and not unkind to the boy about coal baskets and attention, and went back to the anvil as though nothing at all had happened, because as far as the forge was concerned, nothing had.

Arya stood in the yard with her heart going harder than the moment had earned, and understood, standing there, exactly what her own body had just done without asking her permission.

\textit{He didn’t freeze.}

That was the whole of the thought, and it arrived complete, and it disgusted her the moment it finished arriving, because underneath it, plain as the coal smell still hanging in the yard, was the rest of the sentence she had not let herself say: \textit{that’s one more entry on the good side of the page.}

She had a page. She understood that now, standing in a forge yard on an unremarkable afternoon months after she had told a girl of fifteen to hold a retractor steady — she had a page for Gendry, and she had apparently been keeping it the whole while without once deciding to, the same unhurried machinery that had once tallied Ser Meryn Trant and the Tickler and a hundred lesser names now quietly, patiently, without malice and without her leave, tallying the man she loved.

\textit{The machinery’s never once been wrong,} she thought, and felt the whole of the old cold arrive with the thought, familiar as a blade she’d carried too long to feel the weight of anymore. \textit{That was never the problem. The problem is that it doesn’t know how to stop, and I don’t know how to make it, and I am standing in a yard watching a boy I love put out a fire correctly, and some part of me is writing it down.}

\sceneBreak

She told him that evening, because she had decided a long while back that she was not going to be a person who let a true thing rot unsaid simply because saying it was going to cost more than silence.

They were sitting where they always sat, on the low wall behind the forge with the city’s noise going on beyond it, and she waited until the boys had gone home and the coals had banked low, because some things were not owed an audience even of two frightened twelve-year-olds.

“I keep an account on you,” she said. “I want you to know that plainly, because I’ve let it run a long while without ever saying so, and I think you’re owed the saying more than you’re owed my silence about it.”

Gendry did not answer immediately, which she had come, across these months, to understand was simply how he took in anything that mattered.

“The lane,” he said, eventually.

“The lane, and everything since. Today, in the yard — you moved before I’d turned my head, and I stood there and felt some part of me mark it down, same column as the lane’s in, opposite side. I’ve been doing it since the day you told me you were different than the girl who left, and I never once decided to start, and I don’t know how to stop, and I’ve come to believe, sitting with it these last months, that I’m never going to know how to stop, not with you, not with anyone.”

“That’s not the same as what happened in the lane.”

“No. It’s worse.” She said it plainly, because there was no way to say it that made it kinder and she had decided a long while back not to spend effort dressing a true thing up as something softer. “You froze once, Gendry, for reasons that had a face and a street and years of your own childhood behind them, and you’ve not frozen since, not once, not today, not in all the small ordinary tests a person doesn’t know they’re being put through. That’s not a flaw in you. That’s a man who did an entirely human thing one time, under a weight nobody should have to carry, and has been entirely himself ever since.” She looked at the coals rather than at him, because it was easier to finish the thought that way. “The flaw’s in me. I can watch you be exactly that — ordinary, and brave, and kind to a boy who’s frightened himself with a coal basket — and some part of me still opens the ledger. That’s not something you did. That’s something I am, and I don’t think loving you is going to fix it, and I’ve started to believe that staying and pretending otherwise would be a worse thing to do to you than leaving honestly.”

\sceneBreak

He did not argue with her. That was the part she had known would happen and had braced for regardless, because knowing a thing and bearing it turned out, as they so often did, to be two entirely different weights.

“I didn’t think you would,” he said at last, and there was no anger in it, only the particular rueful, tired understanding she had come to rely on in him more than she had ever let herself say aloud. “Fix it, I mean. I don’t think I ever believed that was mine to do. I’ve watched you a long while now, Arya, longer than either of us has let on to the other, and I don’t believe there’s a version of you that puts the ledger down for anybody, myself included. I think I’ve known that since a burnt lane outside the Mud Gate, if I’m honest, and I stayed anyway, because the parts of you that aren’t the ledger are worth considerably more to me than the ledger costs.”

“That’s not a reason to stay with someone.”

“No,” he agreed. “It’s a reason I did, though, for as long as I got to.”

She looked at him properly then, for what she understood, even as she did it, was very likely the last time she was going to let herself look at him this way.

“I’m going,” she said.

“I know.”

“You don’t know where.”

“No,” Gendry said. “I don’t imagine I’m meant to.”

\sceneBreak

She did not tell Sansa the whole of it, because the whole of it belonged to a wall behind a forge and to nobody else, but she went to her two days later, before first light, with Needle wrapped in a length of oiled cloth under her arm.

Sansa was awake already, which did not surprise her; her sister had never slept well in this city and Arya suspected she never would.

“I’m leaving,” Arya said, without preamble, because preambles had never been a thing either of them had much use for with the other, not since the war had burned the childhood cruelty out of both of them and left something plainer standing in its place.

“Where.”

“I don’t know yet. That’s rather the point of it.”

Sansa looked at her for a long moment, and Arya watched her sister do the thing she had watched her do a hundred times this past year — take a piece of news too large to argue with and simply set it down somewhere it could be carried, rather than spend herself fighting a thing that had already been decided.

“Will you write?”

“No,” Arya said, and saw that land, and did not soften it, because she had promised herself on the walk over that she would not dress this morning up as gentler than it was. “I don’t think I will, Sansa. Not because I don’t love you. Because a letter’s a thing with a hand in it and a place it was sent from, and I’ve come to understand, this past while, that those are exactly the two things I need to stop having.”

“Then how will I know you’re alive?”

“You won’t,” Arya said. “I’m sorry. I think that’s the truest answer I’ve got and I’d rather give it to you plain than give you a kinder one that isn’t true.”

Sansa’s eyes were bright, and she did not let them spill over, which Arya thought, watching her, was very likely the single most Sansa thing her sister had ever done in front of her.

Arya held out the cloth-wrapped sword.

“I want you to keep this.”

“It’s yours. Jon gave it to you.”

“I know who gave it to me.” Arya put it into her sister’s hands, because she had decided this on the walk over as well and had found, arriving at the decision, that it did not hurt the way she had expected it to. “I’ve got one working hand and a whole world in front of me that doesn’t know my name, and I’ve decided I’d rather go into it with nothing that says who I was than carry the one thing that would let somebody work it out. Keep it for whoever comes after you, if anyone does. Tell them it was mine, if you like, or don’t — I find I don’t much mind, either way, so long as it’s not in my hand when I go.”

\sceneBreak

She left the city on foot, at the second hour of the morning, through a gate she chose because the man on it was new to the post and had never once seen her face.

He asked her name for the manifest, the way gate-men always asked, because a city kept a record of who went out even when it could not keep one of who came in, and Arya Stark — who had been Arry, and Nan, and a cat with no name of her own on a canal a great many years back, who had worn other people’s faces so many times she had stopped, for a while, being entirely certain which one underneath was hers — found that for the first time in her whole strange life, she did not have an alias ready either.

She had always had one before. That was the thing she noticed, standing in the cold at a gate she’d chosen for its stranger, and found she did not mind noticing.

“Well?” the gate-man said, not unkindly, pen waiting over the page.

Arya Stark looked past him at the road going grey and empty into the last of the dark, and said nothing at all, and after a moment the gate-man shrugged, the way men shrug at travelers too poor or too plain to be worth the trouble of a name, and waved her through, and wrote something in his book that was not her name because she had not given him one to write.

She walked out onto the road with one hand and no ledger anybody else could read, and did not look back at the city, and by the time the sun came fully up there was nothing behind her that knew which way she’d gone.